% ============================================================================
%  PLANTILLA IEEE — Artículo de Conferencia (Conference Paper)
%  Basada en IEEEtran v1.8b+ | Cumple con ieee-latex SKILL.md
%  Uso: copiar este archivo, renombrar, y reemplazar el contenido de ejemplo.
% ============================================================================

\documentclass[conference]{IEEEtran}
\IEEEoverridecommandlockouts
% ↑ Necesario para habilitar \thanks en modo conference.
%   Comentar si no se usa \thanks.

% ====================== PAQUETES ESENCIALES ==================================
% Orden recomendado: encoding → citas → matemáticas → gráficos → tablas → código → hyperref

% --- Encoding y Lenguaje ---
\usepackage[utf8]{inputenc}         % Entrada UTF-8 (acentos, ñ, etc.)
\usepackage[T1]{fontenc}            % Codificación de fuentes para caracteres latinos
% \usepackage[spanish]{babel}       % Descomentar para soporte de español
%                                   % (silabeo, "Resumen" en vez de "Abstract", etc.)

% --- Citas y Bibliografía ---
\usepackage{cite}                   % Ordena y comprime rangos: [1]-[3] en vez de [1],[2],[3]

% --- Matemáticas ---
\usepackage{amsmath,amssymb,amsfonts} % Entornos y símbolos matemáticos extendidos

% --- Gráficos ---
\usepackage{graphicx}               % Inserción de imágenes (SIEMPRE graphicx, nunca psfig)
%                                   % NO especificar driver: \usepackage{graphicx} sin [dvips]
\usepackage[caption=false]{subfig}  % Subfiguras compatibles con IEEE

% --- Tablas ---
\usepackage{array}                  % Columnas mejoradas en tabular
\usepackage{booktabs}               % Líneas profesionales: \toprule, \midrule, \bottomrule
\usepackage{multirow}               % Celdas que abarcan múltiples filas

% --- Código fuente (opcional, para listings) ---
\usepackage{listings}               % Bloques de código con resaltado
\lstset{
  basicstyle=\footnotesize\ttfamily,  % Fuente monoespaciada pequeña
  breaklines=true,                    % Corte automático de líneas largas
  frame=single,                       % Marco alrededor del código
  framesep=6pt,                       % Padding interno entre marco y código
  columns=fullflexible,               % Espaciado flexible (mejor para copiar)
  keepspaces=true,                    % Mantener espacios
  showstringspaces=false,             % No mostrar espacios en strings
  numbers=none,                       % Sin números de línea
  xleftmargin=8pt,                    % Margen izquierdo
  xrightmargin=8pt,                   % Margen derecho
  framexleftmargin=6pt,               % Extensión del marco hacia la izquierda
  framexrightmargin=6pt,              % Extensión del marco hacia la derecha
  aboveskip=10pt,                     % Espacio arriba del bloque completo
  belowskip=6pt,                      % Espacio abajo del bloque completo
  captionpos=t,                       % Caption arriba del listing
  abovecaptionskip=10pt,              % Espacio texto→caption
  belowcaptionskip=15pt,              % FIX: Espacio caption→marco (valor grande)
  language=bash,                      % Lenguaje por defecto
}

% --- Algoritmos (opcional) ---
\usepackage{algorithm}              % Entorno flotante para algoritmos
\usepackage{algorithmic}            % Pseudocódigo (IEEE preferred)

% --- Hyperlinks (SIEMPRE de último) ---
\usepackage{url}                    % Formateo de URLs (\url{})
\usepackage[hidelinks]{hyperref}    % Enlaces clicables sin cajas de color
%                                   % hidelinks = aspecto limpio para impresión

% ====================== FIN DE PAQUETES ======================================

% NOTA: NO cargar geometry, psfig, epsfig, subfigure (viejo), times, ni float.
% NOTA: NO modificar márgenes, fuentes, ni espaciado. IEEEtran lo maneja todo.

% ====================== RUTA DE IMÁGENES (opcional) ==========================
\graphicspath{{figures/}}           % Carpeta de imágenes (evita rutas largas)

% ====================== INICIO DEL DOCUMENTO =================================
\begin{document}

% ---------- TÍTULO -----------------------------------------------------------
% REGLA: NO usar símbolos especiales, notas al pie, ni matemáticas en el título.
\title{Título del Artículo: Descriptivo y Sin Caracteres Especiales}

% ---------- AUTORES ----------------------------------------------------------
% Formato: bloques de nombre + afiliación. De izquierda a derecha, luego abajo.
% El apellido debe ser la última parte del nombre.
\author{
  \IEEEauthorblockN{Primer A. Autor\IEEEauthorrefmark{1},
    Segundo B. Autor\IEEEauthorrefmark{2},
    Tercer C. Autor\IEEEauthorrefmark{1}}
  \IEEEauthorblockA{\IEEEauthorrefmark{1}Programa de Ingeniería,
    Universidad Ejemplo,\\
    Ciudad, País\\
    primer@universidad.edu}
  \IEEEauthorblockA{\IEEEauthorrefmark{2}Departamento de Tecnología,
    Otra Universidad,\\
    Ciudad, País\\
    segundo@otrauniversidad.edu}
}

\maketitle

% ---------- ABSTRACT ----------------------------------------------------------
% REGLA: 100–200 palabras. Sin citas, ecuaciones ni caracteres especiales.
\begin{abstract}
  Escribir aquí un resumen conciso de 100 a 200 palabras.
  El abstract debe indicar el problema abordado, la metodología utilizada,
  los principales resultados obtenidos y las conclusiones más relevantes.
  No incluir citas bibliográficas, ecuaciones matemáticas, ni abreviaturas
  no estándar. Este texto es independiente del resto del artículo y debe
  ser comprensible por sí solo.
\end{abstract}

% ---------- KEYWORDS ----------------------------------------------------------
% REGLA: 5–10 términos. El primero en mayúscula. Orden alfabético.
%        Preferir términos del IEEE Taxonomy.
\begin{IEEEkeywords}
  Keyword uno, keyword dos, keyword tres, keyword cuatro, keyword cinco.
\end{IEEEkeywords}

% ============================================================================
%                            CUERPO DEL ARTÍCULO
% ============================================================================

\section{Introducción}
\label{sec:intro}

Contextualizar el problema, motivar el trabajo, y declarar el objetivo.
Citar fuentes relevantes con \verb|\cite{}|: como se muestra en
\cite{haykin2009neural}, \cite{smith2020example}.

Las abreviaturas se definen en el primer uso en el texto:
Redundant Array of Independent Disks (RAID), incluso si ya aparecen en el
abstract.

% ---------- SECCIÓN DE MARCO TEÓRICO / TRABAJO RELACIONADO -------------------
\section{Trabajo Relacionado}
\label{sec:related}

Revisar brevemente la literatura existente y posicionar el aporte del
artículo actual. Referenciar con \cite{jones2021analysis} y
\cite{doe2022framework}.

% ---------- SECCIÓN DE METODOLOGÍA -------------------------------------------
\section{Metodología}
\label{sec:methodology}

Describir el enfoque, herramientas y procedimiento utilizado.

\subsection{Entorno Experimental}
\label{sec:environment}

Ejemplo de lista (usar solo cuando sea realmente necesario):
\begin{itemize}
  \item Software de virtualización: Oracle VirtualBox 7.0.
  \item Sistema operativo: Ubuntu Server 22.04 LTS.
  \item Herramientas: \texttt{mdadm}, \texttt{fdisk}, \texttt{mkfs.ext4}.
\end{itemize}

\subsection{Procedimiento}
\label{sec:procedure}

Describir los pasos del procedimiento en prosa. Para comandos de terminal,
usar el entorno \texttt{lstlisting}:

\begin{lstlisting}[language=bash, caption={Instalación de mdadm.},
  label={lst:install}]
sudo apt-get update
sudo apt-get install mdadm
\end{lstlisting}

Como se muestra en el Listing~\ref{lst:install}, el primer paso es
instalar la utilidad \texttt{mdadm}.

Para comandos más largos que requieren continuación de línea:

\begin{lstlisting}[language=bash, caption={Creación del arreglo RAID 1.},
  label={lst:raid}]
sudo mdadm --create --verbose /dev/md0 \
    --level=1 \
    --raid-devices=3 /dev/sdb /dev/sdc /dev/sdd
\end{lstlisting}

% ---------- SECCIÓN DE ECUACIONES (ejemplo) ----------------------------------
\section{Modelo Propuesto}
\label{sec:model}

% --- Ecuación simple ---
La capacidad utilizable de un arreglo RAID~1 con $n$ discos de capacidad
$C$ se define como:
%
\begin{equation}
  C_{\text{RAID1}} = C \times 1 = C,
  \label{eq:raid1_capacity}
\end{equation}
%
donde $C$ es la capacidad del disco más pequeño del arreglo.
Nótese que, independientemente de $n$, la capacidad útil se limita a la de
un solo disco, ya que los datos se replican completamente.

% --- Referencia correcta ---
Como se expresa en \eqref{eq:raid1_capacity}, la redundancia tiene un costo
directo en espacio utilizable.

% Al inicio de oración: escribir "Equation" completo
Equation~\eqref{eq:raid1_capacity} muestra que la eficiencia de espacio
es $1/n$.

% --- Ecuación multilínea con align ---
\begin{align}
  \text{Eficiencia} &= \frac{C_{\text{útil}}}{C_{\text{total}}}
    = \frac{C}{n \cdot C}, \label{eq:efficiency} \\
  &= \frac{1}{n}. \label{eq:efficiency_simplified}
\end{align}

% ---------- SECCIÓN DE RESULTADOS --------------------------------------------
\section{Resultados}
\label{sec:results}

Los resultados se resumen en la Table~\ref{tab:comparison} y se ilustran
en Fig.~\ref{fig:example}.

% --- TABLA: caption ARRIBA, \label después de \caption ---
\begin{table}[!t]
  \caption{Comparación de Niveles RAID}
  \label{tab:comparison}
  \centering
  \begin{tabular}{lccc}
    \toprule
    \textbf{Nivel} & \textbf{Discos mín.} & \textbf{Tolerancia}
      & \textbf{Eficiencia} \\
    \midrule
    RAID 0 & 2 & 0 fallos      & 100\% \\
    RAID 1 & 2 & $n-1$ fallos  & $1/n$ \\
    RAID 5 & 3 & 1 fallo       & $(n-1)/n$ \\
    RAID 6 & 4 & 2 fallos      & $(n-2)/n$ \\
    \bottomrule
  \end{tabular}
\end{table}

% --- FIGURA: caption ABAJO, \label después de \caption ---
% REGLA CRÍTICA: \label SIEMPRE después de \caption. Nunca antes.
\begin{figure}[!t]
  \centering
  % Especificar tamaño POR FIGURA (no global).
  \includegraphics[width=\columnwidth]{example-image}
  % ↑ Reemplazar 'example-image' con el nombre real del archivo.
  %   Formatos: PDF o EPS para vectores; PNG/JPG para fotos (300+ DPI).
  \caption{Descripción clara de la figura. Terminar con punto.}
  \label{fig:example}
\end{figure}

% --- SUBFIGURAS ---
\begin{figure}[!t]
  \centering
  \subfloat[Primer caso.]{%
    \includegraphics[width=0.45\columnwidth]{example-image-a}%
    \label{fig:sub_a}}
  \hfil
  \subfloat[Segundo caso.]{%
    \includegraphics[width=0.45\columnwidth]{example-image-b}%
    \label{fig:sub_b}}
  \caption{Comparación visual de dos configuraciones.}
  \label{fig:subfigures}
\end{figure}

% --- FIGURA DOBLE COLUMNA (raro, pero disponible) ---
% \begin{figure*}[!t]
%   \centering
%   \includegraphics[width=\textwidth]{wide_figure.pdf}
%   \caption{Figura que abarca ambas columnas.}
%   \label{fig:wide}
% \end{figure*}

% ---------- SECCIÓN DE DISCUSIÓN ---------------------------------------------
\section{Discusión}
\label{sec:discussion}

Interpretar los resultados en contexto. Comparar con trabajos previos
citados en la Section~\ref{sec:related}. Discutir limitaciones.

% ---------- SECCIÓN DE CONCLUSIONES ------------------------------------------
\section{Conclusiones}
\label{sec:conclusion}

Resumir los hallazgos principales y proponer trabajo futuro.

% ---------- AGRADECIMIENTOS --------------------------------------------------
% Sección sin numerar. "Acknowledgment" sin 'e' (estilo americano IEEE).
% Incluir si el trabajo recibió financiamiento o apoyo.
\section*{Acknowledgment}

Los autores agradecen a [nombre/institución] por [tipo de apoyo].

% ====================== BIBLIOGRAFÍA ========================================
% REGLA: Siempre usar BibTeX con IEEEtran.bst. NUNCA formatear manualmente.
% Compilar: pdflatex → bibtex → pdflatex → pdflatex
\bibliographystyle{IEEEtran}
\bibliography{references}          % Carga references.bib

% ====================== FIN DEL DOCUMENTO ====================================
\end{document}
